% Section 9.1, from lectures/L09/appendix/a_avr_transport.md.

\section{The SPI Peripheral and \code{AvrSpi}}
\label{c9:sec:avr_transport}\label{c9:app:a}
This is the bottom of the stack, and the one place the code finally touches real hardware.
Everything above the seam, the register map, the \code{Interface}, the \code{Uart} driver, was built
and host-tested in Chapters~6 to~8 against a scripted \code{Stub}. Here you implement the
\emph{other} \code{driver::transport::Interface}, the one that drives the ATmega328P's own SPI
peripheral, so that exactly the same driver runs over a real 1\,MHz SPI bus with nothing above the
seam changing.

\code{AvrSpi} is deliberately thin: it configures the SPI master once, then translates the three
seam calls into register accesses. All the protocol knowledge (five-byte transactions, the command
byte, the byte order) already lives in the driver above it.

\subsection{The SPI registers}
\label{c9:sec:spi_registers}
The ATmega328P's SPI is three \code{volatile} memory-mapped registers. This is the \code{volatile}
lesson from modern embedded C++, now on the master side of the wire: the compiler must not cache or
reorder these accesses, because each one is a real bus event.

\begin{booktable}{@{}p{4.4em}p{7.6em}L@{}}
  \thead{Register} & \thead{Bit} & \thead{Meaning} \\
  \midrule
  \code{SPCR} & \code{SPE} & SPI enable. \\
  \code{SPCR} & \code{MSTR} & Master (1) or slave (0). \\
  \code{SPCR} & \code{DORD} & Data order: 0 for most significant bit first, 1 for least. \\
  \code{SPCR} & \code{CPOL} & Clock polarity: 0 means \code{SCK} idles low. \\
  \code{SPCR} & \code{CPHA} & Clock phase: 0 means sample on the leading edge. \\
  \code{SPCR} & \code{SPR1}, \code{SPR0} & \code{SCK} rate select (with \code{SPI2X}). \\
  \code{SPSR} & \code{SPIF} & Transfer complete: set when a byte has finished shifting. \\
  \code{SPSR} & \code{SPI2X} & Double the \code{SCK} rate. \\
  \code{SPDR} & --- & Data register: write to start a transfer, read to get the byte clocked in. \\
\end{booktable}

The transport contract of Part~3 of the protocol (\secref{proto:part3}) pins every one of these. The
ATmega is the \textbf{master}, running \textbf{mode 0} (\code{CPOL = 0}, \code{CPHA = 0}),
\textbf{most significant bit first} (\code{DORD = 0}), with \code{SCK} at
\textbf{f\textsubscript{osc}/16 = 1\,MHz}, the Nano's 16\,MHz clock divided by 16
(\code{SPR1:SPR0 = 01}, \code{SPI2X = 0}).

1\,MHz is not a limit the FPGA imposes: \code{spi_slave} synchronizes \code{sclk} into the 50\,MHz
domain and edge-detects it, so it keeps up with far more than this. It is simply the nearest
prescaler value below the rate at which a breadboard and a passive level shifter stop being
reliable, and it leaves comfortable margin for both.

\subsection{Configuring the master}
\label{c9:sec:configuring}
Configuration is a one-time setup of \code{SPCR}, plus the port pins. On the ATmega328P the SPI pins
are fixed: \code{SCK} is PB5, \code{MOSI} PB3, \code{MISO} PB4, and the hardware \code{SS} PB2. Two
details matter. The first is direction: \code{SCK}, \code{MOSI} and \code{SS} are outputs while
\code{MISO} is an input, and you set that in \code{DDRB}. The second is that \code{SS} (PB2)
\textbf{must be an output} in master mode. If it is left an input and something pulls it low, the
peripheral clears \code{MSTR} and demotes itself to a slave mid-flight. PB2 doubles as the chip
select to the FPGA, so driving it is the transport's job anyway.

\subsection{Releasing the registers on destruction}
\label{c9:sec:releasing}
The constructor is the one place that \emph{sets} bits, so the destructor is the one place that
clears them: \code{AvrSpi} acquires the SPI hardware when it is constructed and \textbf{must release
it when it is destroyed}, leaving the registers exactly as the reset left them. Clear every bit the
constructor set, and nothing else: the \code{DDRB} direction bits for \code{SCK}, \code{MOSI} and
\code{SS}, putting them back to inputs; the \code{SS} bit in \code{PORTB}, which is the chip-select
drive and its pull-up; and all of \code{SPCR}, written to \code{0}, which disables the SPI
peripheral.

\code{MISO} was never configured, so leave it untouched. The destructor is \code{noexcept} and,
because the transport \code{Interface} declares a virtual destructor, is marked \code{override}, so
destroying an \code{AvrSpi} through an \code{Interface&} still runs it. This is plain RAII: whatever
the constructor took, the destructor gives back, so the next owner of those pins --- a second
\code{AvrSpi}, or any other code --- starts from the power-on state rather than a half-configured SPI
master.

\subsection{The three seam methods}
\label{c9:sec:seam_methods}
The seam maps directly onto the peripheral. \code{begin()} drives the chip select low (\code{SS},
PB2), opening a transaction, and \code{end()} releases it high again, closing it. Between them,
\code{transfer(byte)} exchanges one byte in full duplex: write the byte to \code{SPDR}, wait for the
hardware to finish, then read \code{SPDR} for the byte that arrived on \code{MISO}.

That wait is the whole subtlety. A byte takes eight \code{SCK} periods to shift, and \code{SPDR} does
not hold the received byte until the transfer completes, signalled by \code{SPIF} going high:

\begin{cppcode}
SPDR = byte;
while (0U == (SPSR & (1U << SPIF))) {} // Spin until the transfer completes.
return SPDR;                           // Now SPDR holds the received byte.
\end{cppcode}

\noindent Read \code{SPDR} \emph{before} \code{SPIF} is set and you get stale data from the previous
exchange, so the poll is not optional. (Reading \code{SPSR} then \code{SPDR} also clears
\code{SPIF}, arming the next transfer.) This is the same three-method seam the \code{Stub} of
\chapref{8} implemented; the difference is that these bytes are real.

\subsection{Testing it on the host}
\label{c9:sec:host_testing}
The transport touches hardware, but its \emph{logic} is still worth pinning without a bench, so it
is host-tested exactly as the \code{Stub} of \chapref{8} was, over a mocked register file. On the
target, \code{AvrSpi} reaches the registers through a platform header
(\file{arch/avr/hw_platform.hpp}), which forwards to \code{<avr/io.h>} and resolves \code{SPCR},
\code{SPSR}, \code{SPDR}, \code{DDRB} and \code{PORTB} to the real memory-mapped hardware. On the
host, that same platform header selects a provided mock instead, under \code{-DHOST_TEST}, which
backs \code{DDRB}, \code{PORTB} and \code{SPCR} with plain storage and makes \code{SPDR} and
\code{SPSR} proxy objects, so that reading them has an effect the way it does on the part. It records
the bytes written to \code{SPDR} as a \code{MOSI} log and returns a scripted byte on each
\code{SPDR} read from a \code{MISO} queue, and it models the one property that matters here: a
transfer takes time. Writing \code{SPDR} starts it, \code{SPIF} stays \textbf{clear}, and only a
read of \code{SPSR} --- the driver's own poll --- completes the transfer and makes the received byte
readable. Delete the poll and \code{transfer()} hands back a stale byte (on the
mock, the value left from reset, since without the poll no transfer ever completes), and the suite
goes red. It is the direct counterpart of the \code{Stub}.

With that mock, a provided host suite pins the register-level behaviour the bench cannot easily
show: the master configuration bits, the chip-select framing (\code{begin} low, \code{end} high),
the \code{SPIF} poll, and that \code{transfer} returns the byte the mock presented, in order. Only
the register file needs mocking; the transport uses no interrupts and no \code{F_CPU}-derived
timing, so neither \code{<avr/interrupt.h>} nor \code{F_CPU} comes into it. The whole transport is
then proven end to end on the bench in \chapref{10}.

\subsection{What's ahead}
\secref{c9:sec:freestanding} covers the target the transport is built for: a freestanding avr-gcc
build, what it removes, the \file{env.cpp} runtime stubs, and flashing with avrdude. The exercises
in \secref{c9:sec:exercises} implement \code{AvrSpi}, the freestanding bring-up \code{main}, and
UART-based logging.
